\documentclass[11pt,letterpaper]{article}
\usepackage[margin=0.85in,headheight=14pt]{geometry}
\usepackage[T1]{fontenc}
\usepackage{lmodern}
\usepackage{microtype}
\usepackage{amsmath,amssymb}
\usepackage{booktabs,tabularx,array}
\usepackage{enumitem}
\usepackage{xcolor}
\usepackage{fancyhdr}
\usepackage[colorlinks=true,urlcolor=blue!55!black,linkcolor=blue!55!black,citecolor=blue!55!black]{hyperref}
\hypersetup{pdftitle={Reasoning and Recursive Memory: Two-Switch Pilot},pdfauthor={Qining Shen}}

\definecolor{accent}{RGB}{32,67,94}
\pagestyle{fancy}
\fancyhf{}
\fancyhead[L]{\small\textcolor{accent}{Agent Reasoning: Pilot Experiment}}
\fancyhead[R]{\small Qining Shen}
\fancyfoot[C]{\small\thepage}
\renewcommand{\headrulewidth}{0.3pt}
\setlength{\parindent}{0pt}
\setlength{\parskip}{5pt}
\setlist[itemize]{leftmargin=1.4em,itemsep=3pt,topsep=3pt}
\setlist[enumerate]{leftmargin=1.5em,itemsep=4pt,topsep=3pt}
\newcolumntype{L}{>{\raggedright\arraybackslash}X}
\newcommand{\code}[1]{\texttt{#1}}
\newcommand{\unk}{\textsc{Unknown}}
\newcommand{\red}{\textsc{Red}}
\newcommand{\blue}{\textsc{Blue}}

\title{\vspace{-1.7em}\textcolor{accent}{Reasoning and Recursive Memory}\\[3pt]
\large A Diagnostic Report on the Two-Switch Pilot}
\author{Qining Shen}
\date{Experiment: October 7, 2026 \quad Report: October 9, 2026}

\begin{document}
\maketitle
\thispagestyle{fancy}
\vspace{-1.2em}

\textbf{Main finding.} The current model--prompt pipeline does not yet support a reliable
comparison between factual and reasoning-based memory updates. On four basic tasks,
the full-history actor and factual-memory actor each answered one correctly; the
reasoning-memory actor answered none correctly. Nevertheless, the logs separate
representation errors from readout errors: all final factual memories retained the
correct switch values, whereas all final reasoning memories were the same string,
\code{BLUE}. These are diagnostic findings, not evidence for a general ranking of
the two update strategies.

\section{Question and experimental setup}
The motivating question is whether reasoning helps transform an interaction history
into a compact internal state that supports subsequent decisions. Approximate information
states (AIS) motivate separating information preservation from policy performance
\cite{ais}; MEM1 motivates studying recursively updated agent memory \cite{mem1}.
This pilot uses a frozen language model with two update instructions. It does not
reproduce MEM1's reinforcement learning, learn an encoder, or establish AIS conditions.

\textbf{Environment.} Two fixed binary switches, $A,B\in\{0,1\}$, are revealed one at
a time. The correct final action is
\[
 a^\star(A,B)=
 \begin{cases}
 \red, & A=B,\\[-2pt]
 \blue, & A\ne B.
 \end{cases}
\]
Correct choices receive reward $1$; incorrect choices and abstentions receive $0$.
The reported run contains each of the four switch pairs once, with seed $637$ and
\textbf{no distractors or delay turns}. Although the notebook supports longer runs,
those settings were not used in this real-model test. Both memory conditions receive
the same observation sequence for each episode.

\begin{center}
\small
\begin{tabularx}{\linewidth}{@{}lL@{}}
\toprule
Setting & Value used in the reported run\\
\midrule
Model / execution & Qwen2.5-1.5B-Instruct; frozen weights; CPU; greedy decoding\\
Persistent state & 32 model-token slots, padded with $-1$; fixed dimension\\
Generation caps & Update: 64 new tokens; actor: 32; checker: 96\\
Sampling & Four paired episodes; one observation order per switch pair\\
Validation scope & 50 model calls; no generation-cap hits or memory truncation\\
\bottomrule
\end{tabularx}
\end{center}

\textbf{Notation and timing.} We retain $x_t$ for environment state, $o_t$ for observation,
$h_t$ for available history, $z_t$ for internal state, $r_t$ for decision reasoning,
and $a_t$ for action; $r_t$ does not denote reward. During evidence collection,
$a_t=\code{WAIT}$ and $r_t$ is empty. The updater computes
$z_{t+1}=F(z_t,r_t,a_t,o_{t+1})$ using only the old memory and new observation.
The final actor reads the resulting $z_t$. Here $z_t$ encodes text as token IDs;
it is not a learned continuous latent representation.

\newpage
\section{Conditions, measurements, and observed results}
The \textbf{facts} instruction requests observed values in the form
\code{A=...; B=...}, retaining an unknown marker until a value is observed.
The \textbf{reasoning} instruction requests a derived
\code{relation=EQUAL} or \code{relation=DIFFERENT} after both clues are available.
The task prompt explicitly states that either relation suffices to select the door.
These are differences in update instructions, not a claim that the facts model
performs no internal reasoning.

For each final memory we run an actor and a checker using the same model. The
checker has a larger output cap and is evaluated both with $z_t$ alone and with
$z_t+h_t$. The full-history actor is a separate baseline with the same output cap
as the compressed-memory actor. A rule-based \emph{schema decoder} independently
checks explicit switch facts or relation fields. It is not a general semantic judge.
Hidden switch values are available only to the evaluator; the models receive
observations rather than oracle labels. Actual generated-token counts may differ
despite matched caps between conditions.

\begin{center}
\small
\begin{tabularx}{\linewidth}{@{}Lcc@{}}
\toprule
Measurement & Facts & Reasoning\\
\midrule
Final actor correct & $1/4$ (25\%) & $0/4$ (0\%)\\
Checker with $z_t$ correct & $1/4$ (25\%) & $0/4$ (0\%)\\
Checker with $z_t+h_t$ correct & $1/4$ (25\%) & $1/4$ (25\%)\\
Schema decoder returns the correct door & $4/4$ (100\%) & $0/4$ (0\%)\\
Final actor returns \unk & $3/4$ & $4/4$\\
Invalid final-answer format & $0/4$ & $0/4$\\
Mean final memory length & 7 tokens & 1 token\\
Truncated updates & $0/8$ & $0/8$\\
\bottomrule
\end{tabularx}
\end{center}
The \textbf{full-history actor scored $1/4$ (25\%)}. All accuracy denominators include
\unk. In these logs, the actor failures are not caused by the final-answer parser
rejecting malformed output: all actor responses parsed successfully.

\textbf{Per-episode audit.} Let U denote \unk. Episode identifiers below are the
numeric suffixes of \code{ep-0000} through \code{ep-0003}.

\begin{center}
\small
\setlength{\tabcolsep}{4pt}
\begin{tabular}{@{}cccc l c c c@{}}
\toprule
Case & $(A,B)$ & Correct & Full & Facts memory & Actor & Reasoning memory & Actor\\
\midrule
0000 & $(1,0)$ & \blue & \red & \code{A=1; B=0} & U & \code{BLUE} & U\\
0001 & $(0,1)$ & \blue & U & \code{A=0; B=1} & U & \code{BLUE} & U\\
0002 & $(0,0)$ & \red & U & \code{A=0; B=0} & U & \code{BLUE} & U\\
0003 & $(1,1)$ & \red & \red & \code{A=1; B=1} & \red & \code{BLUE} & U\\
\bottomrule
\end{tabular}
\end{center}

\textbf{What the schema score means.} The decoder recognizes switch assignments and
\code{relation=...}, but not a bare door label. Thus the reasoning group's zero schema
score cannot be read as proof that every memory contains no usable task information.
Interpreting \code{BLUE} directly as an answer would match two of the four labels.
However, the identical final state across cases requiring different doors is a real
distinction failure: a decoder given only that state and the fixed task cannot
answer all four correctly. This observation concerns the tested cases, not general
Markov sufficiency.

\newpage
\section{Failure diagnosis and intervention checks}
\textbf{Correct final facts do not imply successful readout.} In case 0000, the
memory was \code{A=1; B=0}, which determines \blue, yet the actor returned
\code{FINAL=UNKNOWN}. The same type of failure occurred in cases 0001 and 0002.
Since both switch values were present and correct, these three final decisions cannot
be explained by loss of the required switch facts. The larger-cap checker did not
resolve them. Prompt interpretation and decision readout therefore require calibration.

\textbf{Intermediate memories can still be wrong.} In the two B-first episodes,
the factual updater initially filled in an unobserved value of A:
\begin{center}
\small
\begin{tabularx}{\linewidth}{@{}clLL@{}}
\toprule
Case & First observation & First facts memory & Final facts memory\\
\midrule
0000 & \code{B=0} & \code{A=0; B=0} & \code{A=1; B=0}\\
0001 & \code{B=1} & \code{A=1; B=1} & \code{A=0; B=1}\\
\bottomrule
\end{tabularx}
\end{center}
Later observations corrected the final states. Consequently, the final 100\% schema
score does not establish that the updater was reliable throughout the trajectory.

\textbf{Reasoning updates did not follow the intended representation.} All four
first-step reasoning memories were \code{RED}, before both switch values were
available; all four final memories were \code{BLUE}. Neither is the requested
relation field. The traces show premature commitments and a constant final output.
They do not identify whether the underlying cause is instruction following, prompt
design, or some other model--interface interaction. Attributing this solely to model
size would exceed the evidence.

\subsection*{Synthetic positive controls}
Two evaluator-constructed memories deliberately contained the wrong relation. The
same actor was run before correction, with an unchanged-memory sham replay, and
after replacing only the relation value with the correct one. E and D below stand
for \code{relation=EQUAL} and \code{relation=DIFFERENT}.

\begin{center}
\small
\begin{tabular}{@{}ccccccc@{}}
\toprule
Case & Correct & Before memory & Before & Sham & Repaired memory & After\\
\midrule
0002 & \red & D & \blue & \blue & E & \red\\
0000 & \blue & E & \red & \red & D & \blue\\
\bottomrule
\end{tabular}
\end{center}
Both constructed cases changed from wrong to correct; sham outputs were unchanged.
This shows sensitivity to the explicit relation field in these examples. The field
occupies 3 tokens for E and 5 for D, so actual input lengths are not identical.
Oracle information was used only for these diagnostic interventions.

\textbf{No natural-error repair was evaluated.} The predefined selector requires a
wrong decision and exactly one incorrect explicit relation field. Actual reasoning
memories were bare \code{BLUE}, so they were ineligible. Zero eligible repairs does
not mean zero natural failures. The two positive controls are excluded from all
main-task accuracy calculations.

\newpage
\section{Interpretation and next experiment}
The run verifies an end-to-end diagnostic implementation and exposes two separate
problems: inadequate task discrimination in the final reasoning memories and
unreliable LLM readout even when factual evidence is complete. The weak full-history
baseline prevents a useful comparative claim about memory strategies.

This pilot used two observations per case, no distractors, and no truncation; it
does not test long-horizon retention or a binding capacity limit. Each pair has one
greedy rollout. Both \blue\ cases were B-first and both \red\ cases A-first, so
observation order and target are confounded. One seed and four cases do not establish
robustness to wording or model choice. Prompt differences are part of the intervention.
AIS performance guarantees require specific sufficiency or approximation conditions
and suitable planning assumptions \cite{ais}; accuracy on these four cases would
not establish them.

\textbf{Recommended next steps (not yet executed).}
\begin{enumerate}
\item \textbf{Calibrate the decision interface.} Test the reader directly on all four
      correct factual memories and both explicit relation values. Simplify input
      presentation and test both observation orders. Require reliable full-history
      performance before using this setup to compare update strategies.
\item \textbf{Validate the update contract.} Use a structured output schema with
      explicit unknown values, flag premature or inconsistent claims, and report
      format adherence separately from task accuracy. Do not repair outputs with
      ground truth during the main experiment. If constrained decoding or retries
      are introduced, document and apply them consistently across conditions.
\item \textbf{Then introduce memory stress.} Add distractors and delays, vary memory
      budgets, and use multiple observation orders and seeds. Continue to distinguish
      information retention, decision readout, and natural-error interventions.
\end{enumerate}

\section*{Provenance and reproduction}
\small
The numerical tables were checked against the saved run's configuration, results,
baselines, model-call logs, memory trajectories, and interventions:
\code{manifest.json}, \code{results.csv}, \code{summary.csv},\\
\code{baselines.csv}, \code{calls.jsonl}, \code{trajectories.jsonl},
\code{interventions.csv}. They refer exclusively to:

\code{real\_model\_smoke\_hf\_20261007T203117938131Z}.

The symbolic demo and the earlier 0.5B test are excluded. The notebook at
\href{https://github.com/russellyaoduphd-cyber/520.637-HW-2/blob/dc91af59480bb7f2299e7ed5c5f3287d22cf4fda/AgentReasoning.ipynb}{commit \code{dc91af5}}
contains the implementation. To reproduce the reported configuration, set
\code{BACKEND="hf"}, \code{DELAY\_PAIRS=((0,0),)},\\
\code{REPEATS\_PER\_SWITCH\_PAIR=1}, \code{MEMORY\_BUDGETS=(32,)},\\
\code{BOOTSTRAP\_SAMPLES=100}, and \code{MAX\_NATURAL\_REPAIRS=1}; retain seed 637
and the generation caps on page 1. The recorded runtime used two CPU threads,
float32 weights, Python 3.12.14, PyTorch \code{2.14.1+cpu}, and Transformers 4.57.1.
Pin \code{MODEL\_REVISION} to:\\
\code{989aa7980e4cf806f80c7fef2b1adb7bc71aa306}.
GPU execution or different software may change outputs despite greedy decoding.

\begin{thebibliography}{9}
\setlength{\itemsep}{1pt}
\bibitem{ais}
J. Subramanian, A. Sinha, R. Seraj, and A. Mahajan.
\emph{Approximate Information State for Approximate Planning and Reinforcement
Learning in Partially Observed Systems}. JMLR 23(12), 1--83, 2022.
\href{https://jmlr.org/papers/v23/20-1165.html}{Paper}.

\bibitem{mem1}
Z. Zhou et al.\ \emph{MEM1: Learning to Synergize Memory and Reasoning for Efficient
Long-Horizon Agents}. arXiv:2506.15841, 2025.
\href{https://arxiv.org/abs/2506.15841}{Paper}.
\end{thebibliography}

\end{document}
